\section{Auditoría de fuentes: una clase, dos rutas desde la neurociencia}

Esta clase no es una conferencia sobre un solo artículo: dos ponentes parten de regiones cerebrales distintas y de niveles de cómputo distintos, y cada uno llega por su lado a la Transformer-like computation. En los primeros 42 minutos, Trenton Bricken parte de la geometría de alta dimensión de la Sparse Distributed Memory (SDM, memoria distribuida dispersa), explica por qué la hypersphere intersection produce pesos exponenciales que aproximan una softmax y propone que el cerebellum (cerebelo) podría implementar un circuito de lectura y escritura similar. En los 39 minutos siguientes, Will Dorrell parte del cognitive map, el hippocampus y la entorhinal cortex, presenta cómo la Tolman--Eichenbaum Machine (TEM) separa la estructura del contenido y cómo la modern Hopfield memory la reformula como attention.

\subsection{Corrección del video y del alcance erróneos de la versión anterior}

El video \texttt{dEFn6nnoC-8} enlazado en la versión anterior es la defensa doctoral de Trenton Bricken, realizada el 31 de marzo de 2025 y publicada en mayo de 2025; no corresponde a esta clase. La grabación oficial correcta es \texttt{L4DC7e6g2iI} de Stanford Online, y la fecha de la clase es el 7 de marzo de 2023. Además, la versión anterior solo cubría al primer ponente y agregaba contenido que no aparece en la clase, como el enrutamiento de MoE, un dashboard de producción, un playbook de gobernanza, incident response y el despliegue intermodal; todos esos párrafos se eliminaron en esta reescritura.

\begin{longtable}{@{}L{0.18\textwidth}L{0.29\textwidth}L{0.41\textwidth}@{}}
\toprule
Nivel de evidencia & Material local & Función didáctica \\
\midrule
Grabación oficial & 1:22:13, dos ponentes, 66 estados didácticos finales & Determina el orden de la clase, las fórmulas, las derivaciones manuscritas, las preguntas y respuestas y las conclusiones \\
Subtítulos oficiales hechos por personas & 1,950 caption y timed transcript & Recuperan la motivación, los límites de las analogías, el alcance de los experimentos y los detalles de neurociencia \\
Artículo de la primera parte & Attention Approximates SDM & Verificar la intersección en alta dimensión, la aproximación exponencial, la SDM continua y el learned coefficient \\
Artículo de la segunda parte & Artículos sobre TEM y el hippocampal Transformer & Verificar la descomposición estructura/contenido, la actualización de la memoria, el grid-like code y su correspondencia con attention \\
Auditoría visual & 40 slide estándar + 26 estados finales manuscritos & Cubrir por completo ambas partes de la clase sin convertir cada trazo escrito en una captura repetida \\
\bottomrule
\end{longtable}

\begin{warningbox}{Límites históricos y de la evidencia}
Esta clase reconstruye la sesión del 2023-03-07. Una correspondencia matemática indica que «cierto cómputo puede implementarse de forma aproximada mediante otra forma»; no demuestra que el cerebro ejecute un Transformer entrenado. Una asignación a regiones cerebrales indica que «el circuito tiene componentes candidatos para implementarlo»; no significa que se haya confirmado experimentalmente. La semejanza entre TEM y el Transformer se da en el nivel de la recuperación de memoria, el estado posicional y la generalización estructural; tampoco es una identidad anatómica neurona por neurona.
\end{warningbox}

\subsection{Pregunta de investigación y ruta de la primera parte}

La portada presenta el trabajo de Bricken y Cengiz Pehlevan; la diapositiva de motivación resume el problema en una frase: la attention del Transformer es muy eficaz, pero la softmax es una heuristic; ¿podría surgir de forma natural a partir de propiedades simples de los vectores de alta dimensión? La diapositiva de Agenda mantiene el orden: primero la intuición visual, después la derivación matemática y luego la discusión de los componentes del Transformer y de su plausibilidad biológica. Este orden importa, porque si la intersección de esferas en alta dimensión se escribe directamente como una función especial, se oculta su intuición como memoria: el «número de direcciones que se traslapan».

\lecturefigure{slide-01-title-attention-sdm.jpg}{Primera parte: Attention Approximates Sparse Distributed Memory}{Diapositiva recuperada de la grabación oficial V001; 00:00:22}

\lecturefigure{slide-02-why-care.jpg}{Por qué vale la pena estudiarlo: de la heuristic softmax a la geometría de la memoria en alta dimensión}{Diapositiva recuperada de la grabación oficial V002; 00:01:20}

\lecturefigure{slide-03-presentation-overview.jpg}{Ruta de la primera parte: SDM, attention, su correspondencia y la biological plausibility}{Diapositiva recuperada de la grabación oficial V003; 00:01:52}

\teachervoice{Bricken dice explícitamente que primero dará una high-level visual intuition y después entrará en las matemáticas, y anima a los estudiantes a interrumpir con preguntas en cualquier momento. Por eso estas notas siguen el mismo orden: primero se dibujan con claridad la escritura, la lectura y la intersección, y después se derivan los pesos exponenciales. No se trata de evitar las fórmulas, sino de que cada símbolo tenga primero un significado operativo.}

\subsection{Las tres preguntas que comparten ambas partes}

Aunque SDM y TEM provienen de tradiciones teóricas distintas, las dos partes indagan tres niveles de relación: en el \textbf{nivel algorítmico}, cómo recupera el modelo los recuerdos pertinentes; en el \textbf{nivel de representación}, cómo se codifican la estructura y el contenido; en el \textbf{nivel de implementación}, qué circuitos neuronales o matrices aprendibles pueden ejecutar el cómputo. El valor de la clase no está en reducir los tres niveles a la frase «el cerebro es un Transformer», sino en comparar cuándo una correspondencia entre niveles es exacta y cuándo es solo una inspiración.

\begin{importantbox}{Escala de evidencia para leer toda la clase}
Distinga, en orden: equivalencia o aproximación de fórmulas, resultados empíricos en modelos entrenados, similitud de representaciones neuronales, posibilidad de implementación en circuitos candidatos y evidencia causal en el cerebro real. Cuanto más adelante en la escala, más exigente es la evidencia requerida. Una fórmula de attention elegante no demuestra por sí sola la función de una región cerebral, y un grid pattern similar tampoco demuestra por sí solo que dos sistemas usen el mismo algoritmo.
\end{importantbox}

\subsection{Resumen de la sección}

La fuente correcta de la clase incluye a dos ponentes y dos rutas complementarias. La primera explica los pesos de attention desde la memoria asociativa de alta dimensión; la segunda explica la generalización estructural y la attention-like retrieval desde los mapas cognitivos y la hippocampal memory. Todo el contenido posterior debe mantener los niveles de evidencia entre el algoritmo, la representación y la implementación biológica.
